\input{preamble.tex}

\title{Отчет по лабораторной работе № 4. \\ Подготовка экспериментального стенда GNS3}
\author{Данила Стариков \\ НПИбд-02-22}
\date{2024}

\begin{document}

\maketitle
\newpage

\tableofcontents

\newpage
\section{Цель работы}
Установка и настройка GNS3 и сопутствующего программного обеспечения.
\newpage
\section{Выполнение работы}
\subsection{Установка GNS3}
\subsubsection{Установка GNS3-all-in-one}
\begin{enumerate}
    \item Через пакетный менеджер установили необходимые пакеты для GNS (Рис. \ref{01}):
        \begin{minted}{bash}
    sudo dnf install gns3-gui gns3-server
        \end{minted}

\begin{center}
    \centering
    \includegraphics[width=0.8\textwidth]{../images/image01.png}
    \captionof{figure}{Установка необходимых пакетов.}
    \label{01}
\end{center}

\end{enumerate}
\subsubsection{Установка GNS3 VM для VirtualBox}
\begin{enumerate}
    \item Далее установили версию виртуальной машины GNS3 VM для VirtualBox с официального репозитория \url{https://github.com/gns3/gns3-gui/releases}, обратив внимание на необходимость совместимости версий. В работе используются версия GNS3 2.2.50:
        \begin{minted}[breaklines,breakbefore=//]{bash}
wget https://github.com/GNS3/gns3-gui/releases/download/v2.2.50/GNS3.VM.VirtualBox.2.2.50.zip
unzip GNS3.VM.VirtualBox.2.2.50.zip
        \end{minted}
    \item Запустили VirtualBox. Выбрали меню {\tt Файл} {\tt Импорт конфигураций...}. Указали месторасположение распакованного образа GNS3 VM.ova. В следующем окне в параметрах импорта выбрали в политике MAC-адреса "Сгенерировать новые MAC-адреса всех сетевых адаптеров". Нажали {\tt Импорт} (Рис. \ref{02}).

\begin{center}
    \centering
    \includegraphics[width=0.8\textwidth]{../images/image02.png}
    \captionof{figure}{GNS3 VM в VirtualBox.}
    \label{02}
\end{center}

    \item Настроили вложенную виртуализацию в VirtualBox через консоль:
        \begin{minted}{bash}
    vboxmanage modifyvm "GNS3 VM" --nested-hw-virt on
        \end{minted}

        Убедились, что опция теперь включена в графическом интерфейсе (Рис. \ref{03}):

\begin{center}
    \centering
    \includegraphics[width=0.8\textwidth]{../images/image03.png}
    \captionof{figure}{Параметры образа виртуальной машины.}
    \label{03}
\end{center}

    \item Настроили сетевой адаптер. Для этого перешли в {\tt Инструменты}, {\tt Сеть} и {\tt Создать} (Рис. \ref{04}).
        В списке появился новый адаптер {\tt vboxnet0}, который затем выбрали в настройки GNS3 VM (Рис. \ref{05}).

\begin{center}
    \centering
    \includegraphics[width=0.8\textwidth]{../images/image04.png}
    \captionof{figure}{Меню настроек графических адаптеров в VirtualBox.}
    \label{04}
\end{center}

\begin{center}
    \centering
    \includegraphics[width=0.8\textwidth]{../images/image05.png}
    \captionof{figure}{Выбор созданного адаптера в настройках образа.}
    \label{05}
\end{center}


\end{enumerate}
\subsubsection{Запуск экземпляра GNS3 в VirtualBox}
\begin{enumerate}
    \item Запустили виртуальную машину GNS3 VM  (Рис. \ref{06}), затем запустили в операционной системе приложение gns3 (Рис. \ref{07}).

\begin{center}
    \centering
    \includegraphics[width=0.8\textwidth]{../images/image06.png}
    \captionof{figure}{Интерфейс GNS3 VM.}
    \label{06}
\end{center}

\begin{center}
    \centering
    \includegraphics[width=0.8\textwidth]{../images/image07.png}
    \captionof{figure}{Интерфейс приложения GNS3.}
    \label{07}
\end{center}

\end{enumerate}
\subsection{Подключение образа оборудования в GNS3}
\subsubsection{Добавление образа маршрутизатора FRR}
\begin{enumerate}
    \item Запустили GNS3 VM, и затем gns3.
    \item В рабочем пространстве GNS3 на левой боковой панели выбрали просмотр маршрутизаторов (Browse Routers), затем нажали на {\tt + New template} (Рис. \ref{09}).

\begin{center}
    \centering
    \includegraphics[width=0.8\textwidth]{../images/image09.png}
    \captionof{figure}{Меню создания нового маршрутизатора.}
    \label{09}
\end{center}

\item В следующем окне выбрали Routers и образ FRR (FRRouting), нажали Install (Рис. \ref{10}).

\begin{center}
    \centering
    \includegraphics[width=0.8\textwidth]{../images/image10.png}
    \captionof{figure}{Список доступных маршрутизаторов.}
    \label{10}
\end{center}

\item В окне с выбором версии для установки выбрали самую актуальную и запустили загрузку (Рис. \ref{11}).

\begin{center}
    \centering
    \includegraphics[width=0.8\textwidth]{../images/image11.png}
    \captionof{figure}{Выбор версии установки FRR.}
    \label{11}
\end{center}

\item После установки открыли настройки маршрутизатора и во вкладке General Settings выбрали {\tt Send the shutdown signal (ACPI)}. Во вкладке "HDD" поставили галочку {\tt Automatically create a config disk on HDD} (Рис. \ref{12}).

\begin{center}
    \centering
    \includegraphics[width=0.8\textwidth]{../images/image12.png}
    \captionof{figure}{Конфигурация маршрутизатора FRR.}
    \label{12}
\end{center}

\end{enumerate}
\subsubsection{Добавление образа маршрутизатора VyOS}
\begin{enumerate}
    \item Для добавления маршрутизатора VyOS предварительно скачали файл образа с репозитория \url{https://github.com/yamadharma/dyos-build/releases}:
        \begin{minted}[breakbefore=//]{bash}
wget https://github.com/yamadharma/dyos-build/releases/download/v1.3.3/vyos-edu.gns3a
wget https://github.com/yamadharma/dyos-build/releases/download/v1.3.3/vyos-1.3.3-amd64.qcow2
        \end{minted}
    \item Скачанный файл добавили в GNS3 через меню {\tt File}, {\tt import appliance} и далее прошли через аналогичные шаги установки маршрутизатора (Рис. \ref{13}).

\begin{center}
    \centering
    \includegraphics[width=0.8\textwidth]{../images/image13.png}
    \captionof{figure}{Выбор установочного файла VyOS.}
    \label{13}
\end{center}

    \item После установки также в настройках во вкладке General Settings выбрали {\tt Send the shutdown signal (ACPI)}. Во вкладке "HDD" поставили галочку {\tt Automatically create a config disk on HDD} (Рис. \ref{14}).

\begin{center}
    \centering
    \includegraphics[width=0.8\textwidth]{../images/image14.png}
    \captionof{figure}{Конфигурация маршрутизатора VyOS.}
    \label{14}
\end{center}

\end{enumerate}

\section{Выводы}
В результате выполнения лабораторной работы установили и настроили GNS3 и сопутствующее программное обеспечение для дальнейшей работы.

\end{document}
